\documentclass[12pt]{article}
\usepackage[a4paper,margin=2.2cm]{geometry}
\usepackage{amsmath,amssymb,amsthm,mathtools}
\usepackage{hyperref}
\usepackage{cleveref}

% 中文（xelatex）
\usepackage{fontspec}
\usepackage{xeCJK}

\usepackage{fontspec}
\usepackage{xeCJK}
\setCJKmainfont{Noto Serif CJK SC}

\title{Least-Squares}
\date{}
\begin{document}
\maketitle

\begin{flushleft}
    \subsubsection{Least-Squares}

    1: 前置: \\
    我们已经从 Ac = u 确定无解 \\
    推出:
    \begin{align}
        \hat{c} \in  \arg \min_{c \in \mathbb{R}^{2}} ||u - Ac||^{2}
    \end{align}
    并得到了 normal equation: \\
    $A^{T}A\hat{c} = A^{T}u$

    现在我们必须要证明:
    $p = A\hat{c}$ 确实是 $Col(A)$ 中距离 u 最近的向量 \\


    2: 已知: \\
    $W  = Col(A),\; p = A\hat{c},\; r = u - p$ \\
    并且$\hat{c}$ 满足 normal equation: \\
    \begin{align}
        A^{T}A\hat{c} = A^{T}u                         \\
        A^{T}(u - A\hat{c}) = 0                        \\
        A^{T}r = 0                                     \\
        r \in NUll(A^{T}) = Col(A)^{\perp} = W^{\perp} \\
        r \perp W
    \end{align}
    所以我们需要证明下面等式: \\
    对于任意另一个系数向量 c 都有: \\
    \begin{align}
        ||u - Ac||^{2} \geq ||u - A\hat{c}||^{2}
    \end{align}
    简称 $X_{1}$


    3: 证明: \\
    任取: \\
    $c \in \mathbb{R}^{2}$ \\
    令: \\
    $q = Ac \in W$, $p = A\hat{c} \in W$ \\
    $p - q \in W$ \\
    我们已知: \\
    $r  =  u - p  \in W^{\perp}$ \\
    因此: \\
    $r \perp p, r \perp q, r \perp (p - q)$ \\

    3.1: $u -q$ 是什么 ?
    $u - q$ 是一个向量 也就是从 $q \to u $ 出发点是 q \\
    那么: \\
    $u - q = (u - p)  + (p - q) $  \\
    我们也可以分成两段路走 先到 p 再到 q \\
    化简下代数式: \\
    \begin{align}
        u - q = (u - p)  + (p - q) \\
        = r + (p - q)              \\
        = r + (A\hat{c} - Ac)      \\
        = r + A(\hat{c} - c)
    \end{align}
    因为 $r \in W^{\perp} \land A(\hat{c} - c) \in W$ \\
    所以我们可以构建一个直角三角形 $\triangle UPQ, \angle P=90^\circ $
    那么存在勾股定理: \\
    \begin{align}
        ||u - q||^{2} = ||r||^{2} + ||A(\hat{c} - c)||^{2}             \\
        ||u - Ac||^{2} = ||u - A\hat{c}||^{2} + ||A(\hat{c} - c)||^{2} \\
        ||A(\hat{c} - c)||^{2} \geq 0                                  \\
        ||u - Ac||^{2} \geq ||u - A\hat{c}||^{2}                       \\
    \end{align}
    所以 $ \hat{c} = \arg \min_{c \in \mathbb{R}^{2}} ||u - Ac||^{2}$ 成立 \\
    $A\hat{c} = proj_{w}(u)$ (投影最短)\\
    证毕:


    4: 投影向量唯一， 但系数未必唯一 \\
    由:
    \begin{align}
        ||u - Ac||^{2} = ||r||^{2} + ||A(\hat{c}-c)||^{2}
    \end{align}
    我们分析上列等式： 它其实是基于毕达哥拉斯定理 \\
    斜边: $||u - Ac||^{2}$ \\
    底边: $||A(\hat{c}-c)||^{2}$ \\
    对边: $||r||^{2}$ \\
    直角三角形 $\triangle UPQ, \angle P=90^\circ $ \\
    问题1: 什么情况下 Q 与 P 最近 ? \\
    斜边 = 对边, 那么几何上最近, 也就是 $||A(\hat{c}-c)||^{2} = 0$

    现在我们设 Q = P, 那么必然存在  $||A(\hat{c}-c)||^{2} = 0 \to A(\hat{c}-c) = 0$ \\
    我们现在分析 式子 $A(\hat{c}-c) = 0 $, $\hat{c}, c$ 是系数向量 \\
    那么 $\hat{c}, c$ 存在哪几种关系呢? \\
    1: 如果 A 的所有列都线性无关并且 $A \neq 0$ 那么有代数式: \\
    \begin{align}
        \sum_{i = 1}^{n} a_{i}(\hat{c} - c) = 0
    \end{align}
    那么代数只有一个解 $\hat{c} - c = 0 \to \hat{c} = c$ 系数唯一, A 的零空间只有 0 \\
    所有的“基准方向”都没有冗余。这意味着没有任意一个方向可以被其他方向替代。

    2: A 存在列线性相关: \\
    \begin{align}
        z = \hat{c} - c     \\
        c = \hat{c} - z     \\
        Az =  0             \\
        Ac = A(\hat{c} - z) \\
        = A\hat{c} - Az     \\
        = A\hat{c} - 0      \\
        = p - 0
    \end{align}
    也就是说 A 零空间中存在可能的无数个 z 向量, z 向量是系数差向量  \\

    总结: 对于客观存在的 $proj_{w}(u)$ 投影向量它是唯一的, 但从 W 零点出发到到达 $proj_{w}(u)$ 的路径组合不是唯一的


    5: 计算 \\
    已知： \\
    \begin{align}
        A = \begin{bmatrix}
                1 & 1 \\
                1 & 1 \\
                0 & 1 \\
            \end{bmatrix}, u = \begin{bmatrix}
                                   4 \\
                                   2 \\
                                   2
                               \end{bmatrix}, \hat{c} = \begin{bmatrix}
                                                            1 \\
                                                            2
                                                        \end{bmatrix}
    \end{align}
    所以: \\
    \begin{align}
        p = A\hat{c} = \begin{bmatrix}
                           1 * 1 + 1 * 2 \\
                           1 * 1 + 1 * 2 \\
                           0 * 1 + 1 * 2 \\
                       \end{bmatrix} = \begin{bmatrix}
                                           3 \\
                                           3 \\
                                           2
                                       \end{bmatrix}                \\
        r = u - p  =  \begin{bmatrix}
                          4 \\
                          2 \\
                          2
                      \end{bmatrix} - \begin{bmatrix}
                                          3 \\
                                          3 \\
                                          2
                                      \end{bmatrix} = \begin{bmatrix}
                                                          1  \\
                                                          -1 \\
                                                          0
                                                      \end{bmatrix} \\
        ||r||^2  = (\sqrt{1^{2} + (-1)^{2} + 0})^{2}  =  2
    \end{align}

    候选一: \\
    \begin{align}
        c^{(1)} = \begin{bmatrix}
                      0 \\
                      2
                  \end{bmatrix}
    \end{align}
    \begin{align}
        q^{(1)} = Ac^{(1)} = \begin{bmatrix}
                                 2 \\
                                 2 \\
                                 2
                             \end{bmatrix}                           \\
        u -q^{(1)} = \begin{bmatrix}
                         4 \\
                         2 \\
                         2
                     \end{bmatrix} - \begin{bmatrix}
                                         2 \\
                                         2 \\
                                         2
                                     \end{bmatrix} =  \begin{bmatrix}
                                                          2 \\
                                                          0 \\
                                                          0
                                                      \end{bmatrix}  \\
        p - q^{(1)} = \begin{bmatrix}
                          3 \\
                          3 \\
                          2
                      \end{bmatrix} -  \begin{bmatrix}
                                           2 \\
                                           2 \\
                                           2
                                       \end{bmatrix} = \begin{bmatrix}
                                                           1 \\
                                                           1 \\
                                                           0
                                                       \end{bmatrix} \\
        ||u - q^{(1)}||^{2} = ||r||^{2}  +  ||p - q^{(1)}||^{2}
        = 4  =  2 + 2
    \end{align}

    候选二: \\
    \begin{align}
        c^{(2)} = \begin{bmatrix}
                      2 \\
                      1
                  \end{bmatrix}
    \end{align}
    \begin{align}
        q^{(2)} = Ac^{(2)} = \begin{bmatrix}
                                 3 \\
                                 3 \\
                                 1
                             \end{bmatrix}                           \\
        u -q^{(2)} = \begin{bmatrix}
                         4 \\
                         2 \\
                         2
                     \end{bmatrix} - \begin{bmatrix}
                                         3 \\
                                         3 \\
                                         1
                                     \end{bmatrix} =  \begin{bmatrix}
                                                          1  \\
                                                          -1 \\
                                                          1
                                                      \end{bmatrix}  \\
        p - q^{(2)} = \begin{bmatrix}
                          3 \\
                          3 \\
                          2
                      \end{bmatrix} -  \begin{bmatrix}
                                           3 \\
                                           3 \\
                                           1
                                       \end{bmatrix} = \begin{bmatrix}
                                                           0 \\
                                                           0 \\
                                                           1
                                                       \end{bmatrix} \\
        ||u - q^{(2)}||^{2} = ||r||^{2}  +  ||p - q^{(2)}||^{2}
        = 3  =  2 + 1
    \end{align}

    $||u - p||^{2} < ||u - q^{(2)}||^{2} <||u - q^{(1)}||^{2} $


    6: homework \\
    一般性证明: \\
    不使用当前具体数值，严格证明： \\
    \begin{align}
        A^{T}(u  - A\hat{c}) = 0
    \end{align}
    推导出: \\
    $\forall c, \; ||u - Ac||^{2} = ||u - A\hat{c}||^{2} + ||A(\hat{c}-c)||^{2} $

    证:  \\
    已知 $ r = u - A\hat{c} \in Col(A)^{\perp}$ \\
    设:
    \begin{align}
        q = Ac \;(\forall a), \; p = A\hat{c}                           \\
        p - q  = (A\hat{c} - Ac = A(\hat{c} - c))\in Col(A)^{\parallel} \\
        u - q  = u - Ac = (u - A\hat{c}) +  (A\hat{c} - Ac )            \\
    \end{align}
    已知 $u - A\hat{c} \in  Col(A)^{\perp} \; \land \; A\hat{c} - Ac \in  Col(A)^{\parallel}$ \\
    我们就可以构建 直角三角形 $\triangle UPQ, \angle P=90^\circ $ \\
    $u - Ac = $ 斜边, $ A\hat{c} - Ac =$ 底边, $u - A\hat{c} = $ 对边 \\
    基于毕达哥拉斯定理 当直角三角形成立存在: \\
    \begin{align}
        ||u - Ac||^{2} = ||u - A\hat{c}||^{2} + ||A(\hat{c}-c)||^{2}
    \end{align}
    证毕：


\end{flushleft}


\end{document}